% chapters/05/section-02.tex
\section{Membaca Masukan yang Berlangsung Terus}
\label{sec:masukan-berlangsung}

Bagian~4.2 telah menjelaskan perbedaan peristiwa dan keadaan tombol. Bagian
ini memakai perbedaan itu untuk menghasilkan gerak yang berlangsung selama
tombol ditahan.

\subsection{Keadaan Seluruh Tombol}

Perintah \code{pygame.key.get_pressed()} mengembalikan keadaan seluruh tombol
pada saat itu, bukan peristiwa. Nilainya dapat ditanyakan untuk tombol mana
pun.

\begin{kodekeluaran}
keys = pygame.key.get_pressed()

if keys[pygame.K_LEFT] or keys[pygame.K_a]:
    direction.x -= 1
\end{kodekeluaran}

Perhatikan susunan kalimatnya: pemeriksaan dilakukan tanpa menunggu peristiwa,
sehingga hasilnya selalu menggambarkan keadaan terkini. Apabila tombol masih
ditahan pada bingkai ini, nilainya benar; apabila sudah dilepas, nilainya
salah. Karena itu cara ini dipakai untuk gerak yang berlangsung terus.

Perhatikan pula bahwa dua tombol dihubungkan dengan \code{or}, sehingga pemain
dapat memakai tombol panah maupun huruf WASD. Kebiasaan menyediakan kedua
pilihan itu penting untuk kenyamanan pemain, terutama pada komputer jinjing
yang tidak selalu memiliki papan tombol angka dan panah yang nyaman dipakai.

\subsection{Menyusun Vektor Arah}

Keempat pemeriksaan tombol disatukan menjadi satu vektor arah. Apabila tidak
ada tombol yang ditekan, vektornya bernilai nol. Apabila satu tombol ditekan,
panjang vektornya satu. Akan tetapi, apabila dua tombol ditekan bersamaan,
misalnya kanan dan bawah, vektornya bernilai (1, 1) dengan panjang
$\sqrt{2} \approx 1{,}414$.

Panjang yang lebih besar itu membawa akibat yang nyata: objek bergerak 41
persen lebih cepat ketika menempuh arah diagonal daripada ketika bergerak
lurus. Perbandingan keduanya diperlihatkan pada
Gambar~\ref{fig:normalisasi-diagonal}.

\begin{figure}[htbp]
\centering
% chapters/05/figures/normalisasi-diagonal.tikz
% Gerak diagonal tanpa dan dengan penormalan vektor arah.
\begin{tikzpicture}[
  panah/.style={-{Stealth[length=2.2mm]}, very thick},
  kisi/.style={draw=ubLine!60, very thin},
  judul/.style={font=\footnotesize\bfseries, ubPrimary},
  label/.style={font=\scriptsize, align=center}
]

% ---------- Tanpa penormalan ----------
\begin{scope}
  \foreach \x in {0,1,2} \draw[kisi] (\x,0) -- (\x,-2);
  \foreach \y in {0,-1,-2} \draw[kisi] (0,\y) -- (2,\y);
  \draw[panah, ubRed] (0,0) -- (1.7,-1.7);
  \node[label, ubRed, anchor=north] at (1.0,-2.15)
    {panjang $\sqrt{2} \approx 1{,}414$};
  \node[judul, anchor=south] at (1.0,0.25) {tanpa penormalan};
  \node[label, anchor=north, text width=42mm] at (1.0,-2.75)
    {kecepatan menjadi $283$ piksel/detik, yaitu $41$ persen lebih cepat};
\end{scope}

% ---------- Dengan penormalan ----------
\begin{scope}[xshift=7.2cm]
  \foreach \x in {0,1,2} \draw[kisi] (\x,0) -- (\x,-2);
  \foreach \y in {0,-1,-2} \draw[kisi] (0,\y) -- (2,\y);
  \draw[panah, ubGreen] (0,0) -- (1.2,-1.2);
  \node[label, ubGreen, anchor=north] at (1.0,-2.15)
    {panjang $1$};
  \node[judul, anchor=south] at (1.0,0.25) {dengan penormalan};
  \node[label, anchor=north, text width=42mm] at (1.0,-2.75)
    {kecepatan tetap $200$ piksel/detik, sama dengan gerak lurus};
\end{scope}
\end{tikzpicture}

\caption{Akibat gerak diagonal tanpa dan dengan penormalan vektor arah}
\label{fig:normalisasi-diagonal}
\end{figure}

Penanganannya sederhana: panjang vektor arah dipaksa menjadi satu sebelum
dikalikan kecepatan.

\begin{kodekeluaran}
if direction.length_squared() > 0:
    velocity = direction.normalize() * speed
else:
    velocity = pygame.Vector2(0, 0)
\end{kodekeluaran}

Pemeriksaan \code{length_squared() > 0} perlu dilakukan lebih dahulu karena
menormalkan vektor bernilai nol tidak bermakna dan dapat menghasilkan nilai
yang tidak diharapkan.

\begin{tipsbox}
Untuk memeriksa apakah vektor bernilai nol, pakailah
\code{length_squared()}, bukan \code{length()}. Keduanya memberi jawaban yang
sama untuk keperluan itu, tetapi \code{length_squared()} lebih murah dihitung
karena tidak menuntut operasi akar. Selisih waktunya tidak terasa pada satu objek,
tetapi menjadi berarti ketika diperiksa untuk ratusan objek setiap bingkai.
\end{tipsbox}

\subsection{Urutan Pembacaan yang Benar}

Vektor arah dibaca \textbf{sekali} pada setiap bingkai, yaitu
pada awal tahap pembaruan, lalu dipakai untuk menghitung kecepatan. Membaca
tombol berulang kali di dalam satu bingkai tidak salah, tetapi tidak ada
gunanya: keadaan tombol tidak berubah di tengah bingkai.

Yang perlu dihindari adalah membaca tombol pada tahap penayangan, misalnya di
dalam fungsi penggambaran. Akibatnya, keadaan yang digambar tidak selalu sama
dengan keadaan yang dihitung, dan jumlah pembacaan bergantung pada berapa kali
fungsi penggambaran dipanggil. Kebiasaan yang benar adalah: masukan dibaca di
tahap masukan, disimpan pada variabel, lalu dipakai oleh tahap-tahap
berikutnya.

\begin{catatabox}
Gerak yang dihasilkan cara ini disebut gerak langsung: objek bergerak dengan
kecepatan penuh begitu tombol ditekan, dan langsung berhenti begitu dilepas.
Permainan yang menghendaki percepatan dan perlambatan---misalnya objek yang
bertambah cepat secara berangsur---menambahkan tahap tersendiri untuk
menyesuaikan kecepatan. Tahap itu dibahas pada Bab~VI.
\end{catatabox}
